% Auto-generated LoRA training examples
% Generated by Scripts/analysis/GenerateResultsFigures.py

% Requires: \usepackage{listings} in preamble

\subsection*{Rounding-Compensation (RC) Example}
Problem: $399 \times 399 = 159201$

\begin{lstlisting}[basicstyle=\small\ttfamily,breaklines=true,literate={×}{{$\times$}}1 {²}{{$^2$}}1]
What is 399 × 399?
Let me round to convenient bases and adjust.
399 is close to 400 (difference: -1).
399 is close to 400 (difference: -1).
Start with 400 × 400 = 160000.
Adjustment for 399: 400 × -1 = -400.
Adjustment for 399: -1 × 400 = -400.
Cross term: -1 × -1 = +1.
Total: 160000 + (-400) + (-400) + (+1) = 159201.
Answer: 159201
\end{lstlisting}

\subsection*{Decomposition (DD) Example}
Problem: $99 \times 40 = 3960$

\begin{lstlisting}[basicstyle=\small\ttfamily,breaklines=true,literate={×}{{$\times$}}1 {²}{{$^2$}}1]
What is 99 × 40?
Let me decompose 99 into 90 + 9.
First compute 90 × 40:
90 × 40 = 3600.
Then compute 9 × 40:
9 × 40 = 360.
Now sum the partial products:
3600 + 360 = 3960.
Answer: 3960
\end{lstlisting}

\subsection*{Ones-Then-Tens (OT) Example}
Problem: $79 \times 78 = 6162$

\begin{lstlisting}[basicstyle=\small\ttfamily,breaklines=true,literate={×}{{$\times$}}1 {²}{{$^2$}}1]
What is 79 × 78?
Let me use column multiplication step by step.
Step 1: Multiply 79 by ones digit 8:
  9 × 8 = 72, write 2, carry 7.
  7 × 8 = 56, plus carry = 63.
  First partial product: 632.
Step 2: Multiply 79 by tens digit 7:
  9 × 7 = 63, write 3, carry 6.
  7 × 7 = 49, plus carry = 55.
  Second partial product: 553 (shifted by 10 = 5530).
Step 3: Add partial products:
  632 + 5530 = 6162.
Answer: 6162
\end{lstlisting}
